\chapter{Conclusion \label{cha:conclusion}}

\section{Summary}

Software modeling and \ac{uml} are important parts of software development. They are included in universities' curriculum not only because they are used in practice, but also because they help students gain a broader perspective on programming and software development - they teach them to plan before writing code and consider the abstract and business aspects of systems.

However, it is common for students and professionals alike to claim knowledge of \ac{uml}, while having only limited or incorrect actual skills in the area.

Traditional education of \ac{uml} faces several challenges such as low amount of feedback, late feedback, or a lack of space for personal attention from teachers. \todo{refine with research findigngs. Students getting overwhelmed by the need to syntax and model at the same time. But that is eh, varime z vody aj ked to je pravda.}

A possible solution to the problem lies in e-learning, which has the potential to improve students' overall learning performance, engagement, and information retention across many domains. It achieves this by allowing for features such as instant feedback, personalized content recommendations, or interactive exercises where students can learn by doing immediately instead of listening to lectures.

% A potential solution to the problem lies in e-learning, which allows not only for instant feedback, but also for personalized content recommendations, and implementation of interactive exercises where students could learn by doing immediately instead of listening to lectures. 

% E-learning has the potential to improve students' engagement, information retention and overall performance across most domains. It uses various practices and designs to drive these outcomes, such as gamification, microlearning, or adaptive learning.

Education of software modeling is no different. Past research has demonstrated promising outcomes by introducing gamified courses, automated diagram assessment tools, exercise recommendation systems and others.

We aimed our research at creating an e-learning course about the \ac{uml} class diagram. Our research goals were twofold. Firstly, the use of e-learning at \ac{fiit} is limited, and we wanted to prepare the ground for further research in the area. Secondly, most research aims at comparing traditional education and e-learning rather than various e-learning approaches. This work contributes to the area by comparing 2 common gamification practices - competition and playfulness. For the experiment, we created 2 variants of the same course applying the 2 designs.

The testing of our course prototypes showed differences in behavioral and emotional engagement related to the 2 designs. However, differences in objective learning outcomes were inconclusive and mostly attributable to differences in starting skill levels of the studied groups. 

Participants in the playful variant exhibited a tendency towards a more explorative approach to the materials, while the competitive variant may promote a more goal-oriented approach directed at measurable outcomes.

Responses to the playful design demonstrated higher interest in encountering more e-learning materials, engagement and overall satisfaction and enthusiasm compared to participants in the competitive design.

Students in the competitive course exhibited more varied responses to the satisfaction metrics, indicating a polarizing effect of competition. 

The overall reception of the courses was positive and interest in encountering more e-learning materials at \ac{fiit} was expressed. These findings suggest that creation of more e-learning content at \ac{fiit} is justified and desirable.


\section{Future Work}

Based on our experiment results and findings from analysis, we recommend conducting further research on the application of e-learning on software modeling - and potentially other areas of software engineering as well. 

Future experiments may extend this work in multiple ways. Tools for enhancing educational processes in software modeling could be integrated into our course - most notably an automated diagram assessment tool, or an exercise generator. The course can be extended to include multiple diagram types, or even new courses aiming other areas of software engineering can be created.

Adaptive and personalized e-learning practices could be explored to address the varying baseline skill levels observed in our experiment. 

Regarding competitive practices, we suggest implementing them in a manner that allows students to make use of the materials without participating in the competition. In-person, group activities might be more suitable for competitive activities.

Playful materials seem accessible to most audiences. Further research may focus on aspects of it not explored in this work, and make an effort to ensure it does not distract students from their learning objectives. 

Experiments could span over longer periods of time so that compound effects of the studied approaches could be demonstrated more clearly.